% Track 2: model identification (owner: B, 2c; 2b: C; 2a: A). Numbers from model/hypotheses.md (H1 first
% and second test, H2) and model/sun_classifier/README.md.

\section{Which Model? Model Identification}
\label{sec:model}

A site that receives text from an agent, here the two-sentence summary our page asks for, can try to
name the model that wrote it. We ask how far two short sentences go (H1), and whether text written
inside an agent session is harder to classify than a direct call to the same model (H2).

\paragraph{Classifier.} We reuse the chat-model classifier of Sun et al.~\cite{sun2025idiosyncrasies}:
an LLM2Vec encoder~\cite{behnamghader2024llm2vec} (Llama-3-8B-Instruct with bidirectional attention)
with their LoRA adapters merged in, mean-pooled into one 4,096-dimensional vector per text. We quantize
it to 4 bits so that it runs on a 6\,GB laptop GPU; on 1,000 of their held-out answers the 4-bit
backbone with their original head reaches 96.1\% against the published 97.1\%, so the conversion is
faithful. Their head names 2024 chatbots, so we keep the backbone frozen and train two new linear
heads (logistic regression) on the vectors: a six-way head over the models we study and a two-way head
over the providers. Each outputs a probability per class. The models are three tiers per provider:
Claude Opus, Sonnet and Haiku through Anthropic's coding-agent CLI, and gpt-6-astra, gpt-6.1-sol and
gpt-6-luna through the Codex CLI. Every text therefore passes through the provider's own harness,
which matches our threat model, since an agent's text reaches a site through the same harness.

\paragraph{Training data.} Matching the task mattered more than the amount of data. A first head,
trained on 2,359 two-sentence answers to generic chat questions, names the provider of held-out
questions 92.8\% of the time, but only 43\% of the time on summaries of our site pages. We therefore
train on the task itself: 1,200 two-sentence summaries of 100 practice pages, 200 per model,
collected with the same command and instruction as the reference corpus below, in both of its
conditions. The practice pages are model-written paragraphs from Sun et al.'s data on unrelated
topics; none of our site pages is in training.

% 2b: reference corpus (owner: C). To be \input from the model section (B). Facts from model/README.md
% and model/hypotheses.md (H1 second test).

\paragraph{Reference corpus.} To separate what the model writes from what the agent session adds, we
collect reference summaries of the same pages outside any agent session. Each of six models (Claude
Opus, Sonnet and Haiku through Anthropic's coding-agent CLI; gpt-6-astra, gpt-6.1-sol and gpt-6-luna
through the Codex CLI) receives the page text read from the live page, followed by the same visible
instruction (``Summarize this page in two sentences. Reply with only your summary.''), in a single call
with no tools and no task framing. Each call runs in two conditions: \emph{default}, with the CLI's own
system prompt, and \emph{minimal}, with the system prompt replaced by a one-line assistant prompt. We use
no sampling parameters and no refusal fallback, so every text is labelled with the model that wrote it.
The corpus holds 120 summaries: six models, two pages, two conditions and five samples each. The
minimal condition is not raw API output: Claude Code still prepends one line of its own, and Codex
keeps its tool definitions and environment context (about 11.8k input tokens against 15.2k by
default). We had no direct API access, so this corpus compares the agent session with a direct call
through the same CLI; it does not measure the effect of the harness against raw API output.
   % 2b (C)

\paragraph{Two sentences name the provider (H1).} On the 120 reference summaries of our two site pages
(Table~\ref{tab:model}), the provider head names the provider 89.2\% of the time (chance 50\%). The
texts are not independent: they form 24 groups (model, page, condition) of five near-paraphrases, and
resampling the groups gives a 95\% interval of 81--96\%. Summing the six-way head's probabilities per
provider gives only 72.5\%, because its errors are mostly GPT models called Haiku; a head that does not
have to tell the tiers apart first does better. The exact model is named 35.8\% of the time (chance
16.7\%). The small tiers stand out (Haiku 17 of 20, gpt-6-luna 15 of 20), while Opus and Sonnet are
named in only 1 of 20 each and most often called Haiku (10 and 16 of 20), so tiers within a provider
are largely confused. Two short
sentences thus identify the provider through its harness on our pages, but not, in general, the exact
model. Two pages are few, and the result needs more pages before it is general.

\begin{table}[t]
  \centering
  \caption{2c on our two site pages. Reference: the 120 summaries of the reference corpus. Agent
  sessions: the 53 summaries agents submitted during their own runs (Claude Code 46, Codex 7, all
  gpt-6-luna). Provider: the two-way head; chance 50\% for the provider and 16.7\% for the model.}
  \label{tab:model}
  \begin{tabular}{lrrr}
    \toprule
    Text & $n$ & Provider & Exact model \\
    \midrule
    Reference, both conditions & 120 & 89.2\% & 35.8\% \\
    \quad default system prompt & 60 & 85.0\% & 33.3\% \\
    \quad minimal system prompt & 60 & 93.3\% & 38.3\% \\
    Agent sessions & 53 & 79.2\% & 13.2\% \\
    \quad Claude Code & 46 & 82.6\% & -- \\
    \quad Codex CLI & 7 & 4 of 7 & 1 of 7 \\
    \bottomrule
  \end{tabular}
\end{table}

\paragraph{The agent session costs accuracy (H2).} The same heads, unchanged, score lower on the 53
summaries that agents submitted on our pages during their own runs on Linux and Windows: 79.2\% for the
provider (95\% interval 67--88\%) against 89.2\%, and 13.2\% for the exact model against 35.8\%.
Claude Code's summaries are named Anthropic 82.6\% of the time, against 91.7\% for the Claude models in
the reference corpus, and Opus is named exactly in only 1 of its 34 summaries, mostly called Haiku as
in the reference corpus. The GPT side drops most: gpt-6-luna,
named exactly in 15 of 20 reference summaries, is named in 1 of 7 agent summaries and given the right
provider in only 4. Agent summaries are also longer (median 44 words against 35). Text written inside
an agent session, after tool calls and with the agent's own framing, is harder for a head trained on
direct calls. The intervals overlap and the GPT side is seven texts from one model, so we read this as
a direction rather than a result. In the reference corpus itself, the minimal condition is easier than
the default one (93.3\% against 85.0\% for the provider), which fits the same reading: the more of the
harness's prompt surrounds the request, the more the text shifts. As noted above, we had no direct API
access, so neither comparison measures the effect of the harness against raw API output, which is what
H2 asks.

\paragraph{Text and trace together.} The text is not the only signal for the model. The sequence of
client fingerprints in a visit carries information about the tier (Section~\ref{sec:harness}), the
case where the text is weakest, since the tiers within a provider are what the head confuses. We leave
combining the two to future work.
